% GENERATED FILE. Edit canonical lessons, then run npm run generate:books.
% canonical-source-hash: fnv1a64:9b4bc0dfc2682dd4
% canonical-lessons: SA-C468-lolupah, SA-C468-pratyupakarah, SA-C468-apakarah, SA-C468-smarayati, SA-C468-krtavedi

\chapter{Gratitude and Greed}
\label{ch:sa-gratitude-and-greed}
\hlchaptermodality{468}

\begin{chapteropening}
\textbf{By the end of this chapter:} \emph{I can say whether someone repays a kindness.}

Complete the last lesson of chapter 468: I can say whether someone repays a kindness.
\end{chapteropening}

\section[\texorpdfstring{\sk{लोलुपः} (lolupaḥ)}{लोलुपः (lolupaḥ)}]{\sk{लोलुपः} (lolupaḥ) — greedy}

\label{lesson:SA-C468-lolupah}



\begin{quote}
Before the new one: say the Sanskrit for flattery, then the Sanskrit for affable.
\end{quote}

\subsection*{You'll want to know: \sk{लोलुपः}}
\textbf{\sk{लोलुपः}} — \emph{lolupaḥ} — \textquotedblleft{}greedy\textquotedblright{}.

\begin{quote}
\sk{लोलुपः} \sk{बालः} \sk{मोदकान्} \sk{खादति}, \sk{मित्राय} \sk{न} \sk{ददाति।}
\end{quote}

\begin{quote}
\emph{lolupaḥ bālaḥ modakān khādati, mitrāya na dadāti.}
\end{quote}

\begin{quote}
The greedy boy eats the sweets and gives his friend none.
\end{quote}

\begin{cousinweb}
\textbf{\sk{लोलुपः}} (\emph{lolupaḥ}) MW builds it on the intensive of the root √lup, \textquotedblleft{}to break\textquotedblright{}, and for the sense \textquotedblleft{}covetous, ardently longing\textquotedblright{} thinks it probably a corruption of \emph{lolubha}, from √lubh, \textquotedblleft{}to covet\textquotedblright{}; its history is not settled (MW).
\end{cousinweb}

\begin{grammarlens}[title={What you've built}]
One more word for describing what you lived and felt.
\end{grammarlens}

\subsection*{Your turn}
\begin{itemize}
\raggedright
  \item \emph{Say it:} \emph{lolupaḥ}
  \item \emph{Say it:} \emph{lolupaḥ}, once more
  \item \emph{Say it:} \emph{lolupaḥ bālaḥ modakān khādati, mitrāya na dadāti.} — The greedy boy eats the sweets and gives his friend none.
  \item \emph{Recall:} say the Sanskrit for flattery, then the Sanskrit for affable, then say \emph{lolupaḥ} again
\end{itemize}

\begin{tcolorbox}[breakable,colback=teal!4,colframe=teal!35!black,title={Before you move on}]
What does \sk{लोलुपः} mean? (\textquotedblleft{}Greedy\textquotedblright{}.) Say it once more.
\end{tcolorbox}
\section[\texorpdfstring{\sk{प्रत्युपकारः} (pratyupakāraḥ)}{प्रत्युपकारः (pratyupakāraḥ)}]{\sk{प्रत्युपकारः} (pratyupakāraḥ) — a favour in return}

\label{lesson:SA-C468-pratyupakarah}



\begin{quote}
Before the new one: say the Sanskrit for affable, then the Sanskrit for greedy.
\end{quote}

\subsection*{You'll want to know: \sk{प्रत्युपकारः}}
\textbf{\sk{प्रत्युपकारः}} — \emph{pratyupakāraḥ} — \textquotedblleft{}a favour in return\textquotedblright{}.

\begin{quote}
\sk{उपकारस्य} \sk{प्रत्युपकारः} \sk{उचितः} \sk{अस्ति।}
\end{quote}

\begin{quote}
\emph{upakārasya pratyupakāraḥ ucitaḥ asti.}
\end{quote}

\begin{quote}
Returning a favour is right.
\end{quote}

A masculine noun.

\begin{cousinweb}
\textbf{\sk{प्रत्युपकारः}} (\emph{pratyupakāraḥ}) is built from \emph{prati}, \textquotedblleft{}back\textquotedblright{}, and \textbf{\sk{उपकारः}}, \textquotedblleft{}a favour\textquotedblright{}; MW gives \textquotedblleft{}returning a service or favour, gratitude\textquotedblright{} (MW).
\end{cousinweb}

\begin{grammarlens}[title={What you've built}]
One more word for describing what you lived and felt.
\end{grammarlens}

\subsection*{Your turn}
\begin{itemize}
\raggedright
  \item \emph{Say it:} \emph{pratyupakāraḥ}
  \item \emph{Say it:} \emph{pratyupakāraḥ}, once more
  \item \emph{Say it:} \emph{upakārasya pratyupakāraḥ ucitaḥ asti.} — Returning a favour is right.
  \item \emph{Recall:} say the Sanskrit for affable, then the Sanskrit for greedy, then say \emph{pratyupakāraḥ} again
\end{itemize}

\begin{tcolorbox}[breakable,colback=teal!4,colframe=teal!35!black,title={Before you move on}]
What does \sk{प्रत्युपकारः} mean? (\textquotedblleft{}A favour in return\textquotedblright{}.) Say it once more.
\end{tcolorbox}
\section[\texorpdfstring{\sk{अपकारः} (apakāraḥ)}{अपकारः (apakāraḥ)}]{\sk{अपकारः} (apakāraḥ) — a wrong; an injury}

\label{lesson:SA-C468-apakarah}



\begin{quote}
Before the new one: say the Sanskrit for greedy, then the Sanskrit for a favour in return.
\end{quote}

\subsection*{You'll want to know: \sk{अपकारः}}
\textbf{\sk{अपकारः}} — \emph{apakāraḥ} — \textquotedblleft{}a wrong; an injury\textquotedblright{}.

\begin{quote}
\sk{सज्जनः} \sk{अपकारं} \sk{क्षमते।}
\end{quote}

\begin{quote}
\emph{sajjanaḥ apakāraṃ kṣamate.}
\end{quote}

\begin{quote}
A good person forgives a wrong.
\end{quote}

A masculine noun.

\begin{cousinweb}
\textbf{\sk{अपकारः}} (\emph{apakāraḥ}) is built from \emph{apa}, \textquotedblleft{}away, off\textquotedblright{}, and \emph{kāra}, \textquotedblleft{}doing\textquotedblright{}, from the root √kṛ of \textbf{\sk{करोति}}, the opposite of \textbf{\sk{उपकारः}}; MW gives \textquotedblleft{}wrong, offence, injury, hurt\textquotedblright{} (MW).
\end{cousinweb}

\begin{grammarlens}[title={What you've built}]
One more word for describing what you lived and felt.
\end{grammarlens}

\subsection*{Your turn}
\begin{itemize}
\raggedright
  \item \emph{Say it:} \emph{apakāraḥ}
  \item \emph{Say it:} \emph{apakāraḥ}, once more
  \item \emph{Say it:} \emph{sajjanaḥ apakāraṃ kṣamate.} — A good person forgives a wrong.
  \item \emph{Recall:} say the Sanskrit for greedy, then the Sanskrit for a favour in return, then say \emph{apakāraḥ} again
\end{itemize}

\begin{tcolorbox}[breakable,colback=teal!4,colframe=teal!35!black,title={Before you move on}]
What does \sk{अपकारः} mean? (\textquotedblleft{}A wrong; an injury\textquotedblright{}.) Say it once more.
\end{tcolorbox}
\section[\texorpdfstring{\sk{स्मारयति} (smārayati)}{स्मारयति (smārayati)}]{\sk{स्मारयति} (smārayati) — reminds}

\label{lesson:SA-C468-smarayati}



\begin{quote}
Before the new one: say the Sanskrit for a favour in return, then the Sanskrit for a wrong done.
\end{quote}

\subsection*{You'll want to know: \sk{स्मारयति}}
\textbf{\sk{स्मारयति}} — \emph{smārayati} — \textquotedblleft{}reminds\textquotedblright{}.

\begin{quote}
\sk{माता} \sk{मां} \sk{प्रतिदिनम्} \sk{औषधं} \sk{स्मारयति।}
\end{quote}

\begin{quote}
\emph{mātā māṃ pratidinam auṣadhaṃ smārayati.}
\end{quote}

\begin{quote}
Mother reminds me about my medicine every day.
\end{quote}

Past (she …): \sk{सा} \sk{मां} \sk{स्मारितवती}, \emph{sā māṃ smāritavatī}.

\begin{cousinweb}
\textbf{\sk{स्मारयति}} (\emph{smārayati}) is built on the causative of the root √smṛ, \textquotedblleft{}to remember\textquotedblright{}, the root of \textbf{\sk{स्मरति}}: to make someone remember; MW's ppp \emph{smārita} is \textquotedblleft{}reminded\textquotedblright{} (MW).
\end{cousinweb}

\begin{grammarlens}[title={What you've built}]
One more word for describing what you lived and felt.
\end{grammarlens}

\subsection*{Your turn}
\begin{itemize}
\raggedright
  \item \emph{Say it:} \emph{smārayati}
  \item \emph{Say it:} \emph{smārayati}, once more
  \item \emph{Say it:} \emph{mātā māṃ pratidinam auṣadhaṃ smārayati.} — Mother reminds me about my medicine every day.
  \item \emph{Recall:} say the Sanskrit for a favour in return, then the Sanskrit for a wrong done, then say \emph{smārayati} again
\end{itemize}

\begin{tcolorbox}[breakable,colback=teal!4,colframe=teal!35!black,title={Before you move on}]
What does \sk{स्मारयति} mean? (\textquotedblleft{}Reminds\textquotedblright{}.) Say it once more.
\end{tcolorbox}
\section[\texorpdfstring{\sk{कृतवेदी} (kṛtavedī)}{कृतवेदी (kṛtavedī)}]{\sk{कृतवेदी} (kṛtavedī) — grateful}

\label{lesson:SA-C468-krtavedi}



\begin{quote}
Before the new one: say the Sanskrit for a wrong done, then the Sanskrit for reminds.
\end{quote}

\subsection*{You'll want to know: \sk{कृतवेदी}}
\textbf{\sk{कृतवेदी}} — \emph{kṛtavedī} — \textquotedblleft{}grateful\textquotedblright{}.

\begin{quote}
\sk{कृतवेदी} \sk{शिष्यः} \sk{गुरुं} \sk{सदा} \sk{स्मरति।}
\end{quote}

\begin{quote}
\emph{kṛtavedī śiṣyaḥ guruṃ sadā smarati.}
\end{quote}

\begin{quote}
A grateful disciple always remembers his teacher.
\end{quote}

\begin{cousinweb}
\textbf{\sk{कृतवेदी}} (\emph{kṛtavedī}) is built from \emph{kṛta}, \textquotedblleft{}what was done\textquotedblright{}, and \emph{vedin}, \textquotedblleft{}knowing\textquotedblright{}, from the root √vid: one who knows what was done for him; MW gives \textquotedblleft{}grateful\textquotedblright{} (MW).
\end{cousinweb}

\begin{grammarlens}[title={What you've built}]
One more word for describing what you lived and felt.
\end{grammarlens}

\subsection*{Your turn}
\begin{itemize}
\raggedright
  \item \emph{Say it:} \emph{kṛtavedī}
  \item \emph{Say it:} \emph{kṛtavedī}, once more
  \item \emph{Say it:} \emph{kṛtavedī śiṣyaḥ guruṃ sadā smarati.} — A grateful disciple always remembers his teacher.
  \item \emph{Recall:} say the Sanskrit for a wrong done, then the Sanskrit for reminds, then say \emph{kṛtavedī} again
\end{itemize}

\begin{tcolorbox}[breakable,colback=teal!4,colframe=teal!35!black,title={Before you move on}]
What does \sk{कृतवेदी} mean? (\textquotedblleft{}Grateful\textquotedblright{}.) Say it once more.
\end{tcolorbox}
